\section*{Capítulo \ref{ch:b3:structural-analysis} --- Análise estrutural}

\begin{solution}{exo:b3:structural-analysis:1}
O \omterm{def:b3:structural-analysis:stereo}{polo} tem rumo $40 + 180 = \ang{220}$ e caimento $90 - 30 = \ang{60}$. O
plano mergulha $90 - 20 = \ang{70}$ em direção a $315 - 180 = \ang{135}$: $135/70$.
\end{solution}

\begin{solution}{exo:b3:structural-analysis:2}
$\theta = \ang{60}$. Ângulos iguais: $10\tan\ang{30} = \qty{5.77}{cm}$;
áreas iguais: $10\sqrt2\sin\ang{30} = \qty{7.07}{cm}$.
\end{solution}

\begin{solution}{exo:b3:structural-analysis:3}
$\vec n_1 = (\sin\ang{30}, 0, \cos\ang{30})$, $\vec n_2 = (-\sin\ang{30}, 0,
\cos\ang{30})$, $\vec n_1\cdot\vec n_2 = -0.25 + 0.75 = 0.5$: os \omterm{def:b3:structural-analysis:stereo}{polos} estão
separados por \ang{60}, o ângulo entre flancos é \ang{120}, uma dobra aberta.
\end{solution}

\begin{solution}{exo:b3:structural-analysis:4}
$\vec d = (\cos\ang{60}, 0, -\sin\ang{60}) = (0.5, 0, -0.866)$ e $\vec u =
0.707\,\vec s + 0.707\,\vec d = (0.354, 0.707, -0.612)$: rumo
$\operatorname{atan2}(0.354, 0.707) = \ang{26.6}$, caimento $\arcsin0.612 =
\ang{37.8}$. O \textit{rake} situa-se entre \ang{-90} (normal) e 0 (sinistral): trata-se
de uma falha normal oblíqua com componente sinistral.
\end{solution}

\begin{solution}{exo:b3:structural-analysis:5}
\begin{gather*}
  \vec n_1 = (0.5567, 0.3214, 0.7660),\qquad \vec n_2 = (-0.4698, -0.1710, 0.8660),\\
  \vec n_1\times\vec n_2 = (0.4093, -0.8420, 0.0558),
\end{gather*}
de comprimento $0.9379$. Seu oposto, $-\vec n_1\times\vec n_2$, aponta para
baixo: rumo \ang{334.1}, caimento $\arcsin(0.0558/0.9379) = \ang{3.4}$. $\vec n_1\cdot\vec
n_2 = 0.3469$, \ang{69.7} entre as normais, ângulo entre flancos \ang{110.3}:
uma dobra aberta.
\end{solution}

\begin{solution}{exo:b3:structural-analysis:6}
Caimento acima de \ang{60} significa $\theta < \ang{30}$. Áreas iguais: $r^2 =
2\sin^2\ang{15} = 0.134$, \qty{13.4}{\percent} da rede; ângulos iguais:
$\tan^2\ang{15} = 0.072$, \qty{7.2}{\percent}. Na esfera, a calota contém
$1 - \cos\ang{30} = 0.134$ do hemisfério: somente a rede de áreas iguais a
mostra com sua parcela verdadeira, por isso os \omterm{def:b3:structural-analysis:stereo}{polos} são contados nela.
\end{solution}

\begin{solution}{exo:b3:structural-analysis:7}
\begin{gather*}
  T_{xx} = \frac{2\sin^2\ang{20}}3 = 0.078,\qquad
  T_{zz} = \frac{2\cos^2\ang{20} + 1}3 = 0.922,\\
  T_{xz} = -\frac{\sin\ang{20}\cos\ang{20} - \sin\ang{20}\cos\ang{20}}3 = 0,
\end{gather*}
e todo termo com $y$ é 0. Os autovalores são $0.922$ (vertical), $0.078$
(leste) e $0$ (norte): o \omterm{def:b3:structural-analysis:pi}{eixo beta} é horizontal, norte--sul. Os três \omterm{def:b3:structural-analysis:stereo}{polos}
estão em um mesmo plano, portanto nenhum tem componente ao longo de sua normal,
e $\vec u^{\mathsf T}T\vec u = 0$ para $\vec u$ ao longo dela.
\end{solution}

\begin{solution}{exo:b3:structural-analysis:8}
$a = 2$, $b = 1.5$, $k = 1/0.5 = 2$: constrição. $Z = 1/(1.8\times1.1) =
0.505$, $a = 1.636$, $b = 2.178$, $k = 0.636/1.178 = 0.54$: achatamento.
\end{solution}

\begin{solution}{exo:b3:structural-analysis:9}
$S_1/S_2 = 2.25$ e $S_1S_2 = 1$: $S_1 = 1.5$, $S_2 = 0.667$; o eixo longo
foi alongado em \qty{50}{\percent}. A cavidade tem os semieixos $1.5\times0.4 =
\qty{0.60}{mm}$ e $0.667\times0.4 = \qty{0.27}{mm}$.
\end{solution}

\begin{solution}{exo:b3:structural-analysis:10}
Ângulos iguais: $1/(2\cos^2\ang{30}) = 0.667$ ao longo do raio e
$\tan\ang{30}/\sin\ang{60} = 0.577/0.866 = 0.667$ ao redor: iguais, conformes.
Áreas iguais: $\mathrm dr/(R\,\mathrm d\theta) = (\sqrt2/2)\cos\ang{30} = 0.612$
e $\sqrt2\sin\ang{30}/\sin\ang{60} = 0.707/0.866 = 0.816$; o produto é
$0.500$, o fator de área do teorema. Um pequeno círculo torna-se uma elipse de
eixos na razão $0.816/0.612 = 1.33$, mais longa na direção em torno do centro.
\end{solution}

\begin{solution}{exo:b3:structural-analysis:11}
A falha é vertical, logo sua normal é horizontal e perpendicular ao
$\sigma_2$ vertical; pela parte (b) do teorema, o deslizamento é perpendicular a
$\sigma_2$, horizontal, para qualquer $\Phi$. Com $\sigma_1 = 1$ ao longo do norte,
$\sigma_3 = 0$ ao longo do leste: $S\vec n = (0, -0.342, 0)$, $\vec n\cdot S\vec n =
0.117$, $\vec\tau = (0, -0.342, 0) - 0.117\,(0.940, -0.342, 0) = (-0.110,
-0.302, 0)$ e $-\vec\tau = (0.110, 0.302, 0)$, em direção a \ang{020}. O bloco
do lado de $\vec n$, o lado leste, move-se para norte-nordeste ao longo da
direção: visto do lado oeste, move-se para a esquerda, uma falha sinistral (\textit{rake}
0). $\sigma_2$ não entrou no cálculo.
\end{solution}

\begin{solution}{exo:b3:structural-analysis:12}
$S = (7.5 - 6.0)/1.0 = \qty{1.5}{km}$; $h = 6.0/1.5 = \qty{4.0}{km}$ abaixo do
horizonte inferior. A seção tinha $20 + 1.5 = \qty{21.5}{km}$ de comprimento e
variou em $-1.5/21.5 = \qty{-7.0}{\percent}$.
\end{solution}

\begin{solution}{pb:b3:structural-analysis:1}
\textbf{1.} O \omterm{def:b3:structural-analysis:stereo}{polo} tem rumo \ang{255} e caimento \ang{77}; $\theta =
\ang{13}$ e $r = 10\sqrt2\sin\ang{6.5} = \qty{1.60}{cm}$ do centro,
em direção a \ang{255}.

\textbf{2.} $\vec n = (\sin\ang{13}\sin\ang{75}, \sin\ang{13}\cos\ang{75},
\cos\ang{13}) = (0.2173, 0.0582, 0.9744)$. Seis vetores unitários quase
paralelos somam quase 6: $|\vec S_E| = 5.9974$.

\textbf{3.} Flanco leste: $\bar\alpha = \operatorname{atan2}(1.2583, 0.2017) =
\ang{80.9}$ e $\bar\delta = \arccos(5.8604/5.9974) = \ang{12.3}$; $\bar R
= 5.9974/6 = 0.99956$ e a dispersão é $\sqrt{2\times0.00044} =
\qty{0.0297}{rad}$, ou \ang{1.7}. Flanco oeste: $\bar\alpha =
\operatorname{atan2}(-1.2445, 0.1886) = \ang{278.6}$ e $\bar\delta =
\arccos(5.8636/5.9972) = \ang{12.1}$; $\bar R = 0.99953$, dispersão \ang{1.8}.

\textbf{4.}
\[
  \vec S_E\times\vec S_W = (0.0774, -14.6714, 0.4883),
\]
de comprimento $14.680$. O seu oposto aponta para baixo: o rumo é $\operatorname{atan2}(-0.0774, 14.6714) = \ang{359.7}$ e o
caimento $\arcsin(0.4883/14.680) = \ang{1.9}$.

\textbf{5.} O cosseno do ângulo entre as normais médias é
$32.835/35.968 = 0.9129$, o produto escalar das somas dividido pelos seus
comprimentos: \ang{24.1} entre as normais, um ângulo interflancos de \ang{155.9}:
uma dobra aberta, próxima de suave.

\textbf{6.}
\begin{gather*}
  \vec m_E = (0.2098, 0.0336, 0.9772),\qquad \vec m_W = (-0.2075, 0.0314, 0.9777),\\
  \vec m_E - \vec m_W = (0.4173, 0.0022, -0.0005):
\end{gather*}
uma normal horizontal voltada para \ang{089.7}. O plano axial é vertical, com direção
\ang{359.7}, praticamente norte--sul: a dobra é reta.

\textbf{7.} O cinturão é o plano cujo \omterm{def:b3:structural-analysis:stereo}{polo} é o \omterm{def:b3:structural-analysis:pi}{eixo beta}: sentido de mergulho
\ang{179.7}, mergulho $90 - 1.9 = \ang{88.1}$, quase o plano vertical leste--oeste
plano. Cada flanco contém o eixo da dobra, logo cada normal é perpendicular a
ele e jaz no plano perpendicular a ele.

\textbf{8.} $0.04400 + 0.00144 + 0.95456 = 1.00000$. $\vec b^{\mathsf T}T\vec
b = 0.00038$, o menor autovalor: $\vec b$ é, a menos do arredondamento, o
autovetor $\vec e_3$, e o cinturão que passa pelos dois \omterm{def:b3:structural-analysis:stereo}{polos} médios é o
cinturão de melhor ajuste entre todos os doze.

\textbf{9.} $\arcsin\sqrt{0.00038} = \arcsin0.0195 = \ang{1.1}$.

\textbf{10.} Ao longo do cinturão, os \omterm{def:b3:structural-analysis:stereo}{polos} distam cerca de \ang{12.2} da vertical,
para cada lado; $\sin^2\ang{12.2} = 0.045$, o quadrado médio de suas
componentes leste--oeste, que é $\lambda_2$.

\textbf{11.} Sim: em média, os \omterm{def:b3:structural-analysis:stereo}{polos} se afastam do cinturão em \ang{1.1}, não
mais do que a dispersão de \ang{1.7} das leituras de um flanco em torno de sua média.

\textbf{12.} $\ang{1.1}/\sqrt{12} = \ang{0.3}$: o caimento é \ang{1.9} $\pm$
\ang{0.3}.

\textbf{13.} Um rake de \ang{-90}: uma falha normal pura, com movimento de mergulho. Anderson:
$\sigma_1$ vertical, $\sigma_3$ horizontal segundo o sentido de mergulho,
\ang{120}--\ang{300}, $\sigma_2$ horizontal segundo a direção,
\ang{030}--\ang{210}.

\textbf{14.} A normal de $120/60$ jaz no plano vertical de rumo
\ang{120}, que contém $\vec e_1$ e $\vec e_3$: $n_2 = 0$. Pela parte (b), o
deslizamento é perpendicular a $\sigma_2$, a direção, para todo $\Phi$: é
um movimento de mergulho. A falha nada informa sobre $\Phi$; na fórmula, tanto $n_2$ quanto
$\vec e_2\cdot\vec w$, logo o numerador e o denominador se anulam.

\textbf{15.} $\vec n = (\sin\ang{55}\sin\ang{150}, \sin\ang{55}\cos\ang{150},
\cos\ang{55}) = (0.4096, -0.7094, 0.5736)$; $n_1 = 0.574$; $n_2 = 0.5\times
0.4096 + 0.866\times(-0.7094) = -0.410$.

\textbf{16.} $\vec e_1\cdot\vec w = 0.198$ e $\vec e_2\cdot\vec w = 0.5\times
0.771 + 0.866\times0.605 = 0.909$, logo $\Phi = -(0.574\times0.198)/(-0.410\times
0.909) = 0.1137/0.3727 = 0.305$.

\textbf{17.} $(0.305 + 0.294 + 0.301 + 0.306)/4 = 0.30$. Os valores de rake foram
arredondados ao grau; um grau de rake desloca $\Phi$ em cerca de 0.02 nestas
falhas, logo o arredondamento de meio grau explica a dispersão.

\textbf{18.} $\sigma_2 - \sigma_3 = 0.3(\sigma_1 - \sigma_3)$: $\sigma_2$ está
mais próxima de $\sigma_3$ do que de $\sigma_1$. Ambas as tensões horizontais estão bem
abaixo da vertical; a região foi distendida sobretudo em direção a \ang{120}, um
pouco segundo \ang{030}.

\textbf{19.} Com $\vec n = (0.750, -0.433, 0.500)$ e o deslizamento $\vec u =
(0.433, -0.250, -0.866)$, $(\vec n - \vec u)/\sqrt2$ tem caimento de \ang{75} para
\ang{300}: P; $(\vec n + \vec u)/\sqrt2$ tem caimento de \ang{15} para \ang{120}:
T. Ambos jazem no mesmo plano vertical que $\sigma_1$ e $\sigma_3$, mas
distam \ang{15} deles, porque a falha forma \ang{30}, e não \ang{45}, com
$\sigma_1$.

\textbf{20.} $\tan a = \tan\ang{12.3}\times\cos(\ang{90} - \ang{80.9}) =
0.2180\times0.9874 = 0.2153$: $a = \ang{12.1}$.

\textbf{21.} $L_0 = 8000/\cos\ang{12} = \qty{8178.7}{m}$, $S =
\qty{178.7}{m}$, uma variação relativa de $\cos\ang{12} - 1 =
\qty{-2.2}{\percent}$.

\textbf{22.} $4000\tan\ang{12} = \qty{850.2}{m}$; $A = \tfrac12\times8000\times
850.2 = \qty{3.40e6}{m^2}$.

\textbf{23.} $h = \num{3.40e6}/178.7 = \qty{19.0}{km}$. Um descolamento na
base da cobertura daria $h$ de algumas centenas de metros e uma área excedente
cerca de cinquenta vezes menor: as dobras não são descoladas ali. Ou o
embasamento dobrou junto com a cobertura, ou os sinclinais afundaram abaixo do
nível original, o que torna $A$ menor; a seção sozinha não permite distinguir.

\textbf{24.} Afastamento horizontal $200/\tan\ang{60} = \qty{115.5}{m}$, deslocamento $200/\sin\ang{60}
= \qty{230.9}{m}$.

\textbf{25.} $\tan a = \tan\ang{60}\cos\ang{30} = 1.500$: $a = \ang{56.3}$, e
o afastamento horizontal aparente é $200/1.5 = \qty{133.3}{m}$. O deslocamento da falha é
no sentido do mergulho, em direção a \ang{120}: a seção ao longo de \ang{120} mostra
comprimentos verdadeiros e é a que deve ser restaurada.

\textbf{26.} A falha corta as dobras, logo é mais jovem e é restaurada primeiro,
movendo o bloco superior de volta ao longo da falha em \qty{230.9}{m}; as dobras são
então desdobradas. As \omterm{def:b3:structural-analysis:restoration}{linhas de fixação} ficam sobre as superfícies axiais dos sinclinais, no
bloco inferior, onde as dobras simétricas não tinham deslizamento entre as camadas.

\textbf{27.} O eixo da dobra de Varden tem caimento
\textbf{\ang[text-series-to-math=true]{1.9}} para o norte (rumo
\ang{359.7}): o \ang{2} embutido na região, recuperado para o
\ang{0.3} que doze leituras permitem.
\end{solution}
